\documentclass[UTF8,fontset=mac,11pt]{ctexart}

\usepackage[a4paper,margin=2.15cm]{geometry}
\usepackage{amsmath}
\usepackage{graphicx}
\usepackage{xcolor}
\usepackage{listings}
\usepackage[hidelinks]{hyperref}

\definecolor{MagneticTBBlue}{HTML}{1455A0}
\definecolor{CodeBackground}{HTML}{F5F7FA}
\definecolor{CodeFrame}{HTML}{D8DEE9}

\hypersetup{
  pdftitle={MagneticTB 安装与快速入门},
  pdfauthor={MagneticTB}
}

\lstdefinelanguage{Wolfram}{
  morekeywords={Needs,PacletInstall,PacletUninstall,PacletFind,Quit,
    FileExistsQ,If,Abort,Total,Dimensions,MatrixForm,Chop,Eigenvalues,
    Thread,Sort,Differences,N,Sqrt,Pi,False},
  sensitive=true,
  morecomment=[s]{(*}{*)},
  morestring=[b]"
}

\lstset{
  language=Wolfram,
  basicstyle=\ttfamily\small,
  keywordstyle=\color{MagneticTBBlue},
  commentstyle=\color{gray},
  stringstyle=\color{teal!70!black},
  backgroundcolor=\color{CodeBackground},
  frame=single,
  rulecolor=\color{CodeFrame},
  framerule=0.4pt,
  framesep=6pt,
  breaklines=true,
  columns=fullflexible,
  keepspaces=true,
  showstringspaces=false,
  aboveskip=0.65em,
  belowskip=0.8em
}

\newcommand{\result}[1]{%
  \par\smallskip
  {\color{gray}\textbf{结果：}}\texttt{#1}
  \par\smallskip
}

\begin{document}

\begin{center}
  {\Huge\bfseries\color{MagneticTBBlue} MagneticTB}\par
  \vspace{0.35em}
  {\LARGE 安装与快速入门}\par
  \vspace{0.7em}
  {\large 从安装到 Hamiltonian、能带和一个简单物性量}
\end{center}

\vspace{0.8em}

这份短指南只保留最常用的主流程，不展示完整 hopping、Wyckoff
或内部数据表。当前安装方式是使用本地 \texttt{.paclet} 文件；安装、更新或卸载后都应重启内核。

\section{安装}

把路径换成实际的 paclet 文件：

\begin{lstlisting}
archive = "/absolute/path/MagneticTB-2.x.y.paclet";
FileExistsQ[archive]
PacletInstall[archive]
Quit[]
\end{lstlisting}

在新内核中确认可以加载：

\begin{lstlisting}
Needs["MagneticTB`"];
PacletFind["MagneticTB"]
\end{lstlisting}

\section{更新}

先确认新 paclet 路径有效，再卸载旧版本并安装新文件：

\begin{lstlisting}
newArchive = "/absolute/path/MagneticTB-2.x.y.paclet";
FileExistsQ[newArchive]

PacletUninstall["MagneticTB"];
PacletInstall[newArchive];
Quit[]
\end{lstlisting}

\section{卸载}

\begin{lstlisting}
PacletUninstall["MagneticTB"]
Quit[]
\end{lstlisting}

\newpage

\section{最小模型：石墨烯}

下面一次完成初始化、构造两带 Hamiltonian，并指定一组数值参数。

\begin{lstlisting}
Needs["MagneticTB`"];

grapheneOperations = msgop[gray[191]];
init[
  lattice -> {{a, 0, 0}, {-a/2, Sqrt[3] a/2, 0}, {0, 0, c}},
  lattpar -> {a -> 1, c -> 3},
  wyckoffposition -> {{{1/3, 2/3, 0}, {0, 0, 0}}},
  symminformation -> grapheneOperations,
  basisFunctions -> {{"pz"}},
  InitialBondShells -> 3,
  GenerateSymmetryGroup -> False
];

h = Total[symham /@ {1, 2, 3}];
parameters = {e1 -> 0.05, r1 -> 0.02, t1 -> 0.5};
\end{lstlisting}

\subsection{查看 Hamiltonian}

不展开冗长的符号表达式，只看维数和 $\Gamma$ 点的数值矩阵：

\begin{lstlisting}
Dimensions[h]
Chop[N[h /. parameters /. {kx -> 0, ky -> 0, kz -> 0}]] // MatrixForm
\end{lstlisting}

\result{\{2, 2\}}

\[
H(\Gamma)=
\begin{pmatrix}
0.17 & 1.50\\
1.50 & 0.17
\end{pmatrix}.
\]

\newpage

\section{画能带}

路径坐标使用倒易点阵的分数坐标。\texttt{bandplot} 在内部把它们转换成
$2\pi\mathbf{k}$ 后计算本征值。

\begin{lstlisting}
path = {
  {{{0, 0, 0}, {0, 1/2, 0}}, {"G", "M"}},
  {{{0, 1/2, 0}, {1/3, 1/3, 0}}, {"M", "K"}},
  {{{1/3, 1/3, 0}, {0, 0, 0}}, {"K", "G"}}
};

bandplot[path, 80, h, parameters, ImageSize -> Large]
\end{lstlisting}

\begin{center}
  \includegraphics[width=0.76\textwidth]{Documentation/figures/quickstart-graphene-band.pdf}
\end{center}

\section{一个简单物性量：K 点直接带隙}

\begin{lstlisting}
kK = 2 Pi {1/3, 1/3, 0};
energiesK = Sort@Eigenvalues@N[
  h /. parameters /. Thread[{kx, ky, kz} -> kK]
];

{energiesK, Chop[Differences[energiesK][[1]]]}
\end{lstlisting}

\result{\{\{-0.01, -0.01\}, 0\}}

因此这组参数下 K 点的两条能带简并，直接带隙为 0。更复杂的 Berry 曲率、
Wilson loop、表面谱等内容留在相应 Mathematica 帮助页中，不塞进这份入门指南。

\end{document}
